% This is the Modeling to the Hodgkin-Huxley neuron model V2.0 book

\chapter{Modeling\label{ch:Modeling}}


\quotationbox{
"What I cannot create, I do not understand"~R.~P.~Feynman}


\quotationbox
{
"Biological laws are much more difficult to find than physical ones"~\cite{VollmerLimitsOfBiology:1995} @1995
}


\section[Methods]{Modeling methods\label{sec:Modeling-Methods}}
%\quotationbox
{"There is a widely accepted distinction between merely modeling a mechanism's behavior and explaining it. \dots Some models are data summaries. Some
models sketch explanations but leave crucial details unspecified or hidden behind
filler terms. Some models are used to conjecture a how-possibly explanation without regard to whether it is a how-actually explanation."\cite{MechanisticModelHodgkinHuxleyCraver:2006}.
}
















\subsection[Disciplinarity]{Disciplinary modeling\label{sec:Modeling-isciplinarity}}
As Feynman R. P~\cite{FeynmanThinking:1980} wrote, "the separation of fields [of science] \dots is merely a human convenience, and an unnatural thing.
Nature is not interested in our separations, and \textit{many of the interesting phenomena bridge the gaps between fields}". "When we go to investigate \textit{we should not pre-decide what it is we are trying to do} except to find out more about it. \dots The first principle is that you must not fool yourself, and you are the easiest person to fool."
Physiology, despite \gls{HH}'s and Feynman's warning, pre-decided that only the science discipline 'electricity' is authorized to describe neurons' operation.
However, electricity has no place for the mass and speed of its charge carriers. After seven decades, it should be realized, admitted, and fixed.



\quotationbox{"Transient electrical signals are particularly
important for carrying \textit{time-sensitive information} rapidly and over long distances. These transient electrical
signals—receptor potentials, synaptic potentials, and
action potentials—are all produced by \textit{temporary
changes} in the electric current into and out of the
cell, changes that drive the electrical potential across
the cell membrane away from its resting value." 
~\cite{JohnstonWuNeurophysiology:1995}
}
